\documentclass[10pt,a4paper]{amsart}
\usepackage[margin=19mm]{geometry}
\usepackage{amsmath,amssymb,mathtools}
\usepackage[scr=rsfs,cal=euler]{mathalfa}
\usepackage{xfrac}
\usepackage[unicode,hidelinks,bookmarks=false]{hyperref}
\newcommand{\ie}{i.e.\ }
\newcommand{\cf}{cf.\ }
\newcommand{\resp}{respectively\ }
\newcommand{\Subsection}{\subsection}
\providecommand{\nobreakdash}{\nobreak\hbox{-}}
\setlength{\emergencystretch}{2em}
\multlinegap0pt
\usepackage{bm}
\usepackage{bbm}
\usepackage{dsfont}
\usepackage{xspace}
\usepackage{cancel}
\usepackage{mathtools}
\usepackage{amsthm}
\makeatletter
\@ifundefined{theorem}{\newtheorem{theorem}{Theorem}}{}
\makeatother
\title[Time reversal of reflected Brownian motion with Poissonian r]{Time reversal of reflected Brownian motion with Poissonian resetting}
\author{Fausto Colantoni, Mirko D'Ovidio, Gianni Pagnini}
\date{Source: arXiv:2505.15639v2 (CC BY 4.0)}
\begin{document}
\maketitle

\noindent\emph{Source and licence.} Time reversal of reflected Brownian motion with Poissonian resetting. Version arXiv:2505.15639v2 (CC BY 4.0), distributed under the Creative Commons Attribution 4.0 International licence, as recorded in the arXiv licence metadata for this exact version and shown on the abstract page https://arxiv.org/abs/2505.15639v2. Full rights notice: licenses/arxiv-2505.15639-CC-BY-4.0.txt.\par\smallskip

\noindent\textit{Editorial scope.} Counted unit: the Theorem whose source label is \texttt{tm:same-local-times} in the article (source0.tex lines 614--621, source label \texttt{tm:same-local-times}), delivered together with the article's own complete original proof of it (source0.tex lines 622--643, one \texttt{proof} environment of 22 lines). No mathematics is written, repaired or re-derived: statement and proof are the article's own text line for line. Required context retained uncounted: tm:local-time-pallino (lines 606--609); symbol:L (lines 1043--1046); drift:joint (lines 1057--1060). Every definition, display and label of those lines is retained unchanged. Omitted from this package and not redistributed: 1--605, 610--613, 644--1042, 1047--1056, 1061--1541. Nothing inside a retained excerpt is omitted or altered: every retained body is delivered exactly as the article's lines are written, so no omission receipt of any kind accompanies this package. Citations: the retained excerpts carry no citation command at all, so no bibliography entry of the article is transcribed and no reference is redistributed. Reference adaptation: none was needed and none was made; every reference inside the delivered unit resolves inside it, so no unresolved reference or citation is produced. Printed slips kept literal: no printed word, symbol, hypothesis or number of the counted statement or of its proof was altered. Layout and class: the article's own class is \texttt{article} with packages including graphicx, amsmath, amssymb, amsfonts, mathrsfs, xcolor, subcaption, hyperref, geometry, amsthm; the wrapper is the lane's pinned \texttt{amsart} editorial wrapper at 10pt on A4, which restates the theorem-like environments the retained text needs and adds the article's own macro definitions that the retained text needs (none), with extra wrapper packages (bm, bbm, dsfont, xspace, cancel, mathtools). The delivered numbering is the unit's own. Assets: no retained span contains a figure environment, a graphics-inclusion command, a PDF-page inclusion, a table or a TikZ picture; no image file is copied into this package, the package declares no asset dependency and no image file is redistributed with it. Licence: version arXiv:2505.15639v2 of this article is distributed under the Creative Commons Attribution 4.0 International licence, as recorded in the arXiv licence metadata for this exact version and shown on the abstract page https://arxiv.org/abs/2505.15639v2. A full rights notice is delivered at licenses/arxiv-2505.15639-CC-BY-4.0.txt. Characters: every retained excerpt is pure ASCII, so the delivered engine, XeLaTeX with its default Latin Modern text font, reports no missing character.
\begin{theorem}
	\label{tm:local-time-pallino}
	The local time at zero of the process $\widetilde{X}$, started at zero, is given by $L^\Psi \circ \widetilde{\gamma}$.
\end{theorem}

	\begin{align}
		\label{symbol:L}
		\int_{0}^\infty e^{-\lambda t} \mathbf{P}_0 (L_t^\Phi \in dx) dt = \frac{\Phi(\lambda)}{\lambda} e^{-x\Phi(\lambda)}.
	\end{align}

	\begin{align}
		\label{drift:joint}
		\int_0^\infty e^{-\lambda t} \mathbf{P}_0(\widetilde{B}_t \in dy, \widetilde{\gamma}_t \in dw) dt = e^{-\sqrt{r} (y-w)} e^{- \sqrt{\lambda + r}(y+w) } dy \, dw.
	\end{align}

\begin{theorem}
	\label{tm:same-local-times}
	Also the inverse of the local time at zero of $\widetilde{X}$, started at zero, is a subordinator with Laplace exponent
	\begin{align*}
		\Phi(\lambda)=\frac{\lambda}{\sqrt{\lambda + r}}, \quad \lambda >0.
	\end{align*}
	When the processes \(X^+\) and \(\widetilde{X}\) start at zero, they share, in distribution, the same local time at zero.
\end{theorem}

\begin{proof}
	By integrating \eqref{drift:joint} in $dy$, we obtain, for $\lambda>0$
	\begin{align*}
		\int_0^\infty e^{-\lambda t} \mathbf{P}_0(\widetilde{\gamma}_t \in du)dt= \frac{1}{\sqrt{\lambda + r} + \sqrt{r}} e^{-w (\sqrt{\lambda + r} - \sqrt{r}) }=\frac{\sqrt{\lambda + r} - \sqrt{r}}{\lambda} e^{-w (\sqrt{\lambda + r} - \sqrt{r}) },
	\end{align*}
	where in the last equality we recognize, by using \eqref{symbol:L}, that $\widetilde{\gamma}$ has the law of the inverse of a subordinator with Laplace exponent $\sqrt{\lambda + r} - \sqrt{r}$. From Theorem \ref{tm:local-time-pallino} we know that the local time of $\widetilde{X}$ is given by $L^\Psi \circ \widetilde{\gamma}$, which are independent, then, for $\lambda>0$,
	\begin{align*}
		\int_0^\infty e^{-\lambda t} \mathbf{P}_0 (L^\Psi \circ \widetilde{\gamma} _ t \in dx) dt &= \int_0^\infty e^{-\lambda t} \int_0^\infty \mathbf{P}_0 (L^\Psi_s \in dx) \mathbf{P}_0(\widetilde{\gamma}_t \in ds) \, dt\\
		&= \int_0^\infty \frac{\sqrt{\lambda + r} - \sqrt{r}}{\lambda} e^{-s (\sqrt{\lambda + r} - \sqrt{r})} \mathbf{P}_0 (L^\Psi_s \in dx) ds\\
		&= \frac{\sqrt{\lambda + r} - \sqrt{r}}{\lambda}  \frac{\Psi(\sqrt{\lambda + r} - \sqrt{r})}{\sqrt{\lambda + r} - \sqrt{r}} e^{-x \Psi(\sqrt{\lambda + r} - \sqrt{r})},
	\end{align*}
	where we have used \eqref{symbol:L}.  We simply observe that
	\begin{align*}
		\Psi(\sqrt{\lambda + r} - \sqrt{r})= \sqrt{\lambda + r} - \sqrt{r} + \sqrt{r} - \frac{r}{\sqrt{\lambda + r} - \sqrt{r} + \sqrt{r}} = \sqrt{\lambda + r} - \frac{r}{\sqrt{\lambda + r}}= \frac{\lambda}{\sqrt{\lambda + r}}.
	\end{align*}
	Then, we conclude from
	\begin{align*}
		\int_0^\infty e^{-\lambda t} \mathbf{P}_0 (L^\Psi \circ \widetilde{\gamma} _ t \in dx) dt&=  \frac{\lambda}{\sqrt{\lambda + r}} \frac{1}{\lambda} e^{-x \frac{\lambda}{\sqrt{\lambda + r}}}\\
		&=  \frac{\Phi(\lambda)}{\lambda} e^{-x\Phi(\lambda)},
	\end{align*}
	that $L^\Psi \circ \widetilde{\gamma}$, which is the local time at zero of $\widetilde{X}$, is the inverse of a subordinator with Laplace exponent $\Phi(\lambda)= \frac{\lambda}{\sqrt{\lambda + r}}$, which holds also for $X^+$.
\end{proof}

\end{document}
